\loai{Tính độ ẩm tuyệt đối}
\begin{vidu}%[1P7-5.2K][Loai1] 
	Buổi sáng nhiệt độ không khí là $20^\circ C$ có độ ẩm tương đối là $70\%$. Cho độ ẩm cực đại ở $20^\circ C$ là $17,3\ g/m^3$. Lượng hơi nước có trong $1\ m^3$ không khí lúc này là
	\choice
	{12,11 g}
	{24,71 g}
	{6,05 g}
	{12,35 g}
	\loigiai{
		Ta có: $a=A.f=12,11\ g/m^3$.
	}
\end{vidu}
\begin{vidu}%[1P7-5.2K][Loai1] 
	Buổi sáng nhiệt độ không khí là $20^\circ C$ có độ ẩm tương đối là $80\%$. Buổi trưa nhiệt độ là $30^\circ C$ có độ ẩm tương đối là $60\%$. Không khí lúc nào chứa nhiều hơi nước hơn? Cho độ ẩm cực đại ở $20^\circ C$ và $30^\circ C$ là $17,3\ g/m^3$ và $30,9\ g/m^3$.
	\choice
	{Buổi sáng}
	{\True Buổi trưa}
	{Đều như nhau}
	{Không xác định được}
	\loigiai{
		\begin{itemize}
			\item Lúc sáng: $a_1=f_1A_1=13,84\ g/m^3$. 
			\item Lúc trưa: $a_2=f_2A_2=18,54\ g/m^3$.
		\end{itemize}
	}
\end{vidu}
\loai{Tính độ ẩm tương đối}
\begin{vidu}%[1P7-5.2K][Loai2] 
	Áp suất hơi nước trong không khí ở $20^\circ C$ là 14,04 mmHg. Cho áp suất hơi bão hòa ở $20^\circ C$ là 17,54 mmHg. Độ ẩm tương đối của không khí trên là
	\choice
	{$60\%$}
	{$70\%$}
	{$80\%$}
	{$85\%$}
	\loigiai{
		Ta có: $f=\dfrac{p}{p_{bh}}=0,8$.
	}
\end{vidu}

\loai{Tính khối lượng nước}
\begin{vidu}%[1P7-5.2K][Loai3]
	Một đám mây có thể tích $100\ km^3$ chứa hơi nước bão hòa ở $20^\circ C$. Vì lí do nào đó, nhiệt độ giảm xuống còn $10^\circ C$ thì lượng nước rơi xuống là bao nhiêu? Cho độ ẩm cực đại ở $20^\circ C$ và $10^\circ C$ là $17,3\ g/m^3$ và $9,4\ g/m^3$.
	\choice
	{$79.10^5$ tấn}
	{\True $7,9.10^5$ tấn}
	{$26,7.10^5$ tấn}
	{$2,67.10^5$ tấn}
	\loigiai{
		\begin{itemize}
			\item Lượng hơi nước có trong đám mây: $m_1=A_1.V=1,73.10^9\ kg$.
			\item Lượng hơi nước có thể chứa ở $10^\circ C$: $m_2=A_2.V=9,4.10^8\ kg$.
			\item Lượng nước mưa: $\Delta m=m_1-m_2=7,9.10^8\ kg=7,9.10^5$ tấn.
		\end{itemize}
	}
\end{vidu}